% =====================================================================
%  RESUMO E ABSTRACT — escreva POR ÚLTIMO (Item 4 das instruções).
%  Pergunta-guia: "O que eu fiz em poucas palavras?"
%  Conteúdo: objetivo (IDÊNTICO ao da Introdução — Dica 7), método e
%  resultados mais relevantes. Um parágrafo, espaçamento simples.
% =====================================================================
\secaosemnumero{RESUMO}
\begin{spacing}{1.0}
\noindent
\escrever{contexto em 1--2 frases: luvas medem movimento, mas não a força nas pontas dos dedos.}
\escrever{objetivo, copiado da Introdução: Este trabalho tem como objetivo ...}
\escrever{método em 2--3 frases: projeto da microestrutura cut-cell em TPU, magnetômetro TLV493D, caracterização em bancada, calibração por regressão e MLP, integração com feedback vibrotátil.}
\escrever{principais resultados com números: RMSE de força em N, faixa útil, latência.}
\escrever{contribuição em 1 frase.}

\vspace{6pt}
\noindent Palavras-chave: Sensor tátil magnético; eFlesh; Manufatura aditiva; Feedback vibrotátil; Aprendizado de máquina.
\end{spacing}

\secaosemnumero{ABSTRACT}
\begin{spacing}{1.0}
\noindent
\begin{otherlanguage}{english}
\escrever{tradução fiel do Resumo para o inglês.}
\end{otherlanguage}

\vspace{6pt}
\noindent Keywords: Magnetic tactile sensor; eFlesh; Additive manufacturing; Vibrotactile feedback; Machine learning.
\end{spacing}
